\onecolumn
\appendices
\section{Original TFM and Study I: Complete Model and PPO Comparison}
\label{app:comparison}
The tables compare the written TFM with Study I. Study II retains these
settings except for the explicit changes in Appendix~\ref{app:v2changes}. ``Not
specified'' is an evidence boundary, not an assertion about missing source code.
All shared changes apply even to our legacy arm. Only reward replacement and
training arrival termination are treatment factors. The exact observation layouts appear in Appendix~\ref{app:observations}.
The accompanying note 25 also indexes source definitions and event ordering.

{\centering\footnotesize
\captionof{table}{Environment, objective, and evaluation comparison. Source locators
refer to printed TFM pages. Rows describe shared changes unless marked as factors.}
\label{tab:modelcomparison}
\begin{tabular}{@{}p{.16\textwidth}p{.37\textwidth}p{.42\textwidth}@{}}
\toprule
Item & Original TFM definition & Study I implementation \\
\midrule
Radio environment & Barcelona, 2.1 GHz, 100 m, 133 sites, nominal 1 m grid;
strongest sector RSS~\cite[pp.~6--7, Sec.~3.1, Table~2]{thesis}
& Received map retained; 87 station operator retained. No ray tracing regeneration;
exact source version and sentinel remain unverified. \\
Lookup and geometry & 3D building propagation geometry; lookup and integer
encoding not specified~\cite[pp.~6--7, Sec.~3.1]{thesis}
& Native coordinate spacing, nearest lookup, float radio arithmetic; $-128$
treated as absent path. No collision model. \\
Motion & Scalar speed/selected heading update; five acceleration and eight
heading bins~\cite[p.~9, Eq.~(2); pp.~12, 15]{thesis}
& Inertial velocity, continuous radial/lateral acceleration, norm bound 5 m/s$^2$,
trapezoidal position update, hard speed cap 25 m/s. \\
Time and speed & 1 ms in prose, 0.1 s in table; 2000 timesteps; 100 m/s
penalty threshold~\cite[pp.~14--16, Tables~3--4]{thesis}
& 1 s steps, 200 s deadline. New time discretization and speed constraint. \\
Arrival (factor) & 10 m tolerance, stopping intention, no arrival termination
\cite[p.~9, Sec.~3.2; p.~12, Sec.~5.2; p.~15, Table~3]{thesis}
& Distance $\leq10$ m and speed $\leq2$ m/s. Training termination varied;
every evaluation ends at first safe arrival. \\
Failure handling & Energy termination; boundary/RSS/speed penalties
\cite[p.~12, Sec.~5.2, Eq.~(14)]{thesis}
& Energy exhaustion, boundary crossing, five consecutive RSS outages terminate.
Physical failure precedes arrival; a valid deadline arrival succeeds. \\
Observations & Position, serving BS, RBG, remaining energy, buffer
\cite[p.~12, Sec.~5.2]{thesis}
& 42 scaled features including goal bearing/distance, velocity, time,
consumed energy, prior arrival, phases, serving SINR, candidates, and masks. \\
Network actions & Station and RBG requests, A3 inequality
\cite[p.~9, Eq.~(3); p.~12, Sec.~5.2]{thesis}
& Stay plus top four RSS stations; $>3$ dB advantage and $>-96$ dBm;
five second decisions plus observed outage requests; no time to trigger counters. \\
Resource allocation & Policy requests RBG; availability fallback; random
occupancy of at least half the groups~\cite[pp.~9, 12, 14--15]{thesis}
& RBG is not learned. First free group, exactly six occupied groups per station
in a deterministic pattern; global label phase is not independent load variation. \\
Rate and interference & Bandwidth plus log of SNR; neighboring RSS sum labeled
uplink~\cite[p.~10, Eqs.~(7)--(9)]{thesis}
& Linear downlink cochannel interference, SINR, $B\log_2(1+\mathrm{SINR})$.
Retain 1.44 MHz and $-103.41$ dBm combined noise; change link interpretation. \\
Queue and traffic & Poisson formulation, ratio delay, 160 KB, 4800 bit handover
overhead; conflicting data packet sizes~\cite[pp.~9--10, Eqs.~(4)--(6); pp.~14--15]{thesis}
& Deterministic 200 kbit/s, decimal 1.28 Mbit cap, same handover overhead.
Postservice backlog/rate censored at 200 s; record overflow without terminating. \\
Energy & Mass/efficiency/lift-to-drag expression; capacity labeled 1000 kW
\cite[p.~10, Eq.~(10); p.~15, Table~3]{thesis}
& Consumed energy integral of $100+0.4v^2$ W; 100 kJ budget. New uncalibrated
proxy, not a conversion of original capacity. \\
Reward (factor) & Positive reciprocals, binary $+12$ progress, penalties
\cite[p.~12, Eqs.~(13)--(14); p.~15, Table~3]{thesis}
& Legacy structure uses new common physical inputs and no soft speed penalty.
Replacement adds first arrival/terminal failure payments, time cost, bounded
radio costs, and terminal corrected potential shaping. \\
Deterministic reference & Constant speed, closest discrete heading, viability
counters~\cite[pp.~13--14, Sec.~5.4]{thesis}
& Continuous desired velocity tracking with braking; strongest admissible
station and shared masks. A new reference, not a reproduction of that greedy policy. \\
Splits and comparison & PPO/greedy and weight comparisons; separate validation
and test sets not specified~\cite[pp.~15--16, Secs.~6.2--6.3]{thesis}
& 4096 training, 64 validation, 200 standard test, 200 longer test routes;
five seeds per treatment, frozen budget. Same map, not a geographic holdout. \\
Reported outcomes & SNR, outage, interference, remaining energy, handovers,
example paths~\cite[pp.~16--20, Sec.~7]{thesis}
& Mission and recorded feasible success first; secondary metrics conditional
on success; every seed, paired bootstrap and failure counts retained. \\
\bottomrule
\end{tabular}
}

For completeness, our reset sets velocity, queue, consumed energy, counters,
and prior arrival to zero, associates with the strongest station, and allocates
its first free group. Each step applies the requested admissible handover before
motion, evaluates radio after motion, updates queue and energy, checks failure
and arrival, and then computes reward. The TFM does not fully specify this
event ordering~\cite[pp.~9--12, Secs.~3.2--5.2]{thesis}.
The direct energy and packet overflow reward terms removed during the original
development~\cite[pp.~20--21, Sec.~7.3]{thesis} are not treated as part of its
final reward. We also do not claim to have reproduced its mass or efficiency
model, random traffic, counter based handovers, or learned resource allocation.

\clearpage
{\centering\footnotesize
\captionof{table}{PPO comparison. The first eleven rows cover every parameter in the
original Table 4~\cite[p.~16]{thesis}. Further implementation details are not
fully specified in its Secs. 5.3 and 6.2~\cite[p.~13; pp.~15--16]{thesis}.}
\label{tab:ppocomparison}
\begin{tabular}{@{}p{.25\textwidth}p{.20\textwidth}p{.50\textwidth}@{}}
\toprule
Parameter & Original TFM & Study I implementation \\
\midrule
Episodes & 300 & Fixed interactions, variable number of episodes \\
Episode timesteps & 2000 & 200, with explicit early terminals \\
Rollout steps per environment & 8 & 128 \\
Learning rate & $3\cdot10^{-5}$ & $3\cdot10^{-4}$, constant \\
Batch size & 32 & 512 \\
Discount & 0.99 & 0.99, retained \\
GAE lambda & 0.95 & 0.95, retained \\
Policy clip & 0.2 & 0.2, retained \\
Value coefficient & 0.5 & 0.5 applied to half MSE; total coefficient 0.25 MSE \\
Entropy coefficient & 0.01 & 0.005 \\
Target KL & 0.03 & No threshold; log diagnostic KL only \\
\midrule
Implementation & Stable Baselines & Native PyTorch hybrid PPO \\
Parallel lanes and update epochs & Not specified & 64 lanes; 8192 samples per rollout; four epochs \\
Interaction budget & Episodes/timesteps listed & 524,288 per policy, 64 rollout updates \\
Networks & Not specified & Separate actor/critic, each two 64 unit tanh layers \\
Policy heads & Discrete action vector listed & 2D Gaussian plus masked five choice categorical \\
Initialization and variance & Not specified & Orthogonal weights; log standard deviation initialized $-0.5$,
clamped to $[-2.3,0.5]$ \\
Optimizer and gradient cap & Adam in pseudocode & Adam, epsilon $10^{-5}$, gradient norm cap 0.5 \\
Likelihood and entropy details & Not specified & Latent Gaussian likelihood after stored action sampling;
sum latent Gaussian and categorical entropy \\
Reward and advantage scaling & Not specified & Running discounted return variance; no reward centering;
reward clip $[-10,10]$; minibatch advantage normalization \\
Critic and bootstrap details & Not fully specified & Unclipped half MSE; bootstrap only across nonterminal
rollout ends, not across mission/deadline terminals \\
Evaluation and checkpoint rule & Not specified & Gaussian mean and categorical mode; final checkpoint only;
common raw evaluation reward for every treatment \\
\bottomrule
\end{tabular}
\vspace{1.5em}
}

\bigskip
\section{Every Executed Change from Study I to Study II}
\label{app:v2changes}
This inventory compares final frozen implementations, including choices that
were retained. Original locators refer to the written TFM, not recovered code.
The inherited serialized field \texttt{max\_consecutive\_outage: 5} is unused by
Study II: its executed step fails on the first violation. Remaining queue at
arrival and instantaneous switching are retained limitations.

{\footnotesize\setlength{\tabcolsep}{3pt}
\begin{longtable}{@{}p{0.13\textwidth}p{0.23\textwidth}p{0.31\textwidth}p{0.28\textwidth}@{}}
\caption{Final Study II differences and original source relationships.}\label{tab:v2changes}\\
\toprule Item & v1 executed behavior & Final v2 executed behavior & Original source relationship\\\midrule
\endfirsthead
\multicolumn{4}{l}{\textit{Continued from previous page}}\\\toprule
Item & v1 executed behavior & Final v2 executed behavior & Original source relationship\\\midrule
\endhead
\bottomrule\endfoot
Minimum RSS & RSS $\leq$ -96 dBm counts as outage. & RSS < -96 dBm counts as violation; equality is allowed. & Follow C2's $\geq$ inequality \cite[p. 11, Eq. (12)]{thesis} and the numeric minimum \cite[p. 15, Table 3]{thesis}.\\
Initial RSS & Strongest station selected at reset; no separate persistent initial violation flag. & Check the initial serving RSS and preserve any violation until failure. & Include the initial modeled sample in C2 \cite[p. 11]{thesis}.\\
Failure rule & Five consecutive sampled outages terminate a flight. & Any single sampled RSS violation terminates a failed joint mission. & Enforce the written C2 at modeled samples \cite[p. 11]{thesis}, rather than reproduce the source's soft reward penalty \cite[p. 12, Eq. (14)]{thesis}.\\
Buffer overflow & Count lost bits, clip queue, continue flight; disqualify only the subsequent feasible label. & Any overflow ends failure; clipping cannot erase it. & Enforce C4 at the sampled queue update \cite[p. 11, Eq. (12)]{thesis}. Source removal of a direct dump penalty is described on pp. 20--21, Sec. 7.3.\\
Joint success & Arrival is primary; recorded feasibility permits up to 5\% RSS outage and zero drops. & Arrival plus no prior sampled RSS violation, no overflow, and no physical failure. No additional permissive feasibility label. & Combine source communication constraints \cite[p. 11]{thesis} with explicit destination completion.\\
Event precedence & Some physical failures override arrival; overflow does not. & RSS, overflow, boundary and energy failure override simultaneous arrival. & New explicit ordering; original source ordering is unspecified \cite[p. 12, Sec. 5.2]{thesis}.\\
Outstanding queue at arrival & Allowed. & Still allowed and reported. & The source does not specify an empty queue arrival test \cite[pp. 9--12, Secs. 3.2--5.2]{thesis}. We do not claim all generated data has been delivered.\\
Delay cost & 0.05 * min(D, 1) inside fixed reward. & 0.35 * 10D / (1 + 10D) in full arm. No flat reward plateau at one second. & One minus original reciprocal utility \cite[p. 12, Eq. (13)]{thesis} with equal configuration values \cite[p. 15, Table 3]{thesis}. Smooth saturation remains.\\
Interference cost & Absent from fixed reward, although measured. & Restore 0.30 * 100000I / (1 + 100000I). & Same reciprocal utility scales and weights \cite[pp. 12, 15]{thesis}. I remains our downlink proxy.\\
Handover cost & 0.05 * H. & 0.35 * 100H / (1 + 100H). & Restore normalized source utility \cite[p. 12, Eq. (13); p. 15, Table 3]{thesis}.\\
Positive per step radio utility & Present in reimplemented legacy reward, removed in v1 fixed reward. & Subtract communication cost, so simply staying active does not earn that utility bonus. & Changes the positive reciprocal reward structure \cite[p. 12, Eqs. (13)--(14)]{thesis} while preserving the three normalized tradeoff components.\\
Learning comparison & Four reward/termination arms; direct acceleration learning. & Two arms: arrival base only versus arrival base minus all three radio costs. Both use strict failure, termination and the same new controller architecture. & A new controlled comparison; neither arm is the source's trained policy.\\
Network decision timing & Every five seconds, with observed outage exception. & Every modeled second. & Source describes choices each environment step \cite[pp. 9, 12, Secs. 3.3.1, 5.2]{thesis}; a five second interval was our prior assumption.\\
Network action & Stay plus four strongest candidate stations: five categorical choices. & Stay plus five candidate station slots by 12 groups: 61 choices. & Restores resource choice from the source's BS/RBG action \cite[p. 12, Sec. 5.2]{thesis}.\\
Station candidate selection & Serving plus strongest four, stable RSS ordering. & Retained. Duplicate serving station slots are masked. & Restricted candidate interface is ours; source lists full BS/RBG choices \cite[p. 12]{thesis}.\\
A3 condition & Candidate RSS > serving RSS + 3 dB. & Retained; prospective filter cannot bypass it. & Follow the inequality, not the reversed nearby prose \cite[p. 9, Eq. (3); p. 15, Table 3]{thesis}.\\
Resource allocation & Deterministic first free RBG at reset and handover. & That initialization remains; the controller can choose any currently free RBG, including within the same station. & Source allows RBG requests with fallback \cite[p. 9, Sec. 3.3.1]{thesis}. The exact masking implementation is ours.\\
Repeating a resource choice & No learned resource choice. & Current serving station/group pair remains a valid pair action; explicit stay also remains. & Avoid a new artificial alternating action mask; this is an implementation decision.\\
Resource change count & Not independently reported. & Count every station/group change, including same station RBG changes. & Diagnostic addition. Same station resource changes are not counted as handovers.\\
Switching interruption & Not modeled. & Still not modeled; no extra RBG switching cost. & Source does not provide a calibrated interruption duration \cite[p. 9, Sec. 3.3.1; p. 15, Table 3]{thesis}. This remains a limitation, not a verified original behavior.\\
Continuous action & Two direct radial/lateral acceleration commands. & Projected reference command plus two bounded learned residuals of scale 0.5, followed by physical projection. & New controller prior; differs from original discrete acceleration and heading \cite[p. 12, Sec. 5.2; p. 15, Table 3]{thesis}.\\
Safety filter & None. & Predict the next sample; repair unsafe network choice if possible, otherwise try local motion/network alternatives; still fail if unresolved. & New model based constraint handling component; original source describes reward penalties \cite[p. 12, Eq. (14)]{thesis}.\\
Prospective radio access & Learned policy observes current candidates. & PPO still observes current options. Its shared filter and model based references query the known map at predicted positions. & New information/control architecture, not reproduced source logic.\\
Observation & 42 features. & 156 features, detailed below. & Original source has six listed state elements \cite[p. 12, Sec. 5.2]{thesis}. Both are reimplementations with repaired state.\\
PPO network head & Five categorical logits. & 61 categorical logits; continuous and value heads retained. & Original architecture details are not sufficiently specified \cite[p. 13, Sec. 5.3]{thesis}.\\
Training route lengths & 200--1000 m. & 200--1800 m; longer tests lie within this range. & New study design; original route distribution and independent splits are not documented \cite[pp. 15--16, Secs. 6.2--6.3]{thesis}.\\
Final route identities & Seeds 42012 and 42013. & Fresh seeds 52012 and 52013. Exact pairs are disjoint from v1 and other v2 splits. & Our explicit split protocol; same map remains shared.\\
Training seeds & 1101--1105. & 2101--2105. & New repeated experiment; original seed list is unspecified \cite[pp. 15--16]{thesis}.\\
Final budget & 524,288 interactions each; 20 policies. & Same interactions per policy; 10 policies. & Our PPO budget remains different from 300 episodes and 2,000 maximum steps in the source \cite[p. 16, Table 4]{thesis}.\\
Deterministic references & Destination braking and strongest admissible RSS. & Four nominal control laws with the same filter: Goal RSS, Goal radio, one step joint search, and three step joint search. & New references, not a recreation of the thesis's discrete constant velocity greedy controller \cite[pp. 13--14, Sec. 5.4]{thesis}.\\
Secondary comparison & Success conditioned summaries and selected matched comparisons. & Explicit intersections of successful route/seed pairs for radio comparisons, with counts and crossed bootstrap intervals. & Source does not report a matching conditional joint success protocol \cite[pp. 16--20, Sec. 7]{thesis}.\\
\end{longtable}}
\clearpage
\section{Exact Observation and Numerical Interfaces}
\label{app:observations}
All indices are zero based. Study I has 42 features; Study II has 156. Both
replace the original listed position, station, resource, energy and queue vector
\cite[p.~12, Sec.~5.2]{thesis}. Scales are fixed; no fitted observation normalizer
is used. Study II removes the prior arrival flag, consecutive outage count and
five second phase because arrival and RSS failure now terminate immediately.

{\footnotesize\setlength{\tabcolsep}{3pt}
\begin{longtable}{@{}p{0.13\textwidth}p{0.82\textwidth}@{}}
\caption{Study I observation indices and scales.}\label{tab:obs1}\\
\toprule Indices & Features and scaling\\\midrule
\endfirsthead
\multicolumn{2}{l}{\textit{Continued from previous page}}\\\toprule
Indices & Features and scaling\\\midrule
\endhead
\bottomrule\endfoot
0--1 & Position divided by (5000, 3500)\\
2--3 & Unit vector toward goal; (1, 0) fallback at effectively zero distance\\
4 & Goal distance / 1000 m\\
5--6 & Velocity projected onto radial and lateral directions / 25 m/s\\
7 & Speed / 25 m/s\\
8 & Elapsed steps / 200; determines remaining time\\
9 & Queue bits / 1,280,000\\
10 & Consumed energy / 100,000 J, rather than remaining energy\\
11 & Whether first safe arrival already occurred\\
12--13 & Resource label phase / 11 and allocated RBG / 11\\
14 & (step mod 5) / 5\\
15 & Serving SINR in dB / 30\\
16 & Consecutive outage samples / 5\\
17--21 & Five candidate RSS values, each transformed as (RSS + 70) / 60\\
22--31 & Candidate station x/y offsets from UAV, both divided by 5000 m\\
32--36 & Candidate indices within selected operator / 86, not global station IDs or one hot vectors\\
37--41 & Five Boolean admissibility flags represented as float features\\
\end{longtable}}

{\footnotesize\setlength{\tabcolsep}{3pt}
\begin{longtable}{@{}p{0.13\textwidth}p{0.82\textwidth}@{}}
\caption{Study II observation indices and scales.}\label{tab:obs2}\\
\toprule Zero based indices & Definition\\\midrule
\endfirsthead
\multicolumn{2}{l}{\textit{Continued from previous page}}\\\toprule
Zero based indices & Definition\\\midrule
\endhead
\bottomrule\endfoot
0--1 & Position / map dimensions\\
2--3 & Destination direction unit vector\\
4 & Destination distance / 1,000 m\\
5--6 & Radial and lateral velocity / 25 m/s\\
7 & Speed / 25 m/s\\
8 & Elapsed decision count / 200\\
9 & Queue / 1,280,000 bits\\
10 & Consumed energy / 100,000 J\\
11--12 & Background phase / 11 and serving RBG / 11\\
13 & Serving SINR in dB / 30\\
14--18 & Candidate RSS, transformed as (RSS + 70) / 60\\
19--28 & Candidate station x/y offsets / 5,000 m\\
29--33 & Candidate station index / 86\\
34--94 & log(1 + option\_capacity / 200000) for all 61 network choices\\
95--155 & Corresponding Boolean action masks as float features\\
\end{longtable}}

For station index $k$, episode label phase $z$, and group $g$, background
occupancy is $\ind[(g-(7k+z))\bmod12<6]$. The reset allocator chooses
$(7k+z+6)\bmod12$. Phase $z$ rotates labels globally; it does not yield independent
load fields. Each vector lane is an independent single UAV mission. The source
describes random occupancy of at least half the groups
\cite[p.~14, Sec.~6.1; p.~15, Table~3]{thesis}.

Both studies use postservice backlog divided by service rate, censored at 200 s,
as delay. A zero service rate is handled numerically rather than interpreted
as a measured packet latency. The shared queue generates offered bits and any
handover overhead, serves the queue at the modeled capacity, records excess
above its capacity and clips for storage. Study II additionally fails on that
excess. Stored $-128$ RSS supplies zero power; SINR reporting has a $10^{-15}$
logarithm floor. These declarations complement the original queue and rate
expressions~\cite[pp.~9--10, Eqs.~(4)--(9)]{thesis}.

Study I evaluates every checkpoint using the common fixed raw reward by default;
its stored return is an undiscounted evaluation sum, not its PPO normalized
training return. Success and physical metrics define the reported comparisons.
In continuing training, first arrival pays its bonus once; a later failure
cannot pay the prearrival failure penalty. Study II always terminates on joint
success or first failure and scores the proposed residual/network action before
environment projection and safety filtering. Neither study has a separate soft
speed penalty, since the shared dynamics impose the cap.

\clearpage
\section{Development Evidence and Artifact Map}
\label{app:development}
Development results below use only 64 standard and 32 longer validation routes.
Each row uses one seed per arm and is exploratory; it does not estimate final
seed uncertainty. All six Study II pilot checkpoints and source snapshots are
retained locally. Study I retains six pilots with seeds 41 and 42, including
the validation curve in Fig.~\ref{fig:pilot}. No pilot checkpoint initializes a
final policy.

{\footnotesize\setlength{\tabcolsep}{3pt}
\begin{longtable}{@{}p{0.34\textwidth}p{0.21\textwidth}p{0.2\textwidth}p{0.2\textwidth}@{}}
\caption{Study II pilot development; percentages are standard / longer joint success.}\label{tab:development}\\
\toprule Design & Pilot seed and budget per arm & Arrival arm standard / longer joint success & Full arm standard / longer joint success\\\midrule
\endfirsthead
\multicolumn{4}{l}{\textit{Continued from previous page}}\\\toprule
Design & Pilot seed and budget per arm & Arrival arm standard / longer joint success & Full arm standard / longer joint success\\\midrule
\endhead
\bottomrule\endfoot
Direct motion, no filter, short training lengths & 51; 524,288 & 89.1\% / 59.4\% & 92.2\% / 34.4\%\\
Direct motion, safety filter, short training lengths & 52; 1,048,576 & 92.2\% / 78.1\% & 98.4\% / 75.0\%\\
Residual motion, safety filter, 200--1800 m training lengths & 53; 524,288 & 96.9\% / 84.4\% & 98.4\% / 87.5\%\\
\end{longtable}}

The first strict nominal RSS controller completed 95.3\% / 71.9\% of validation
routes; prospective radio selection completed 96.9\% / 90.6\%. A local joint
search initially completed 95.3\% / 81.3\% and sometimes preferred stationary
low cost states. A preference for immediate progress, with the shared filter,
raised its observed validation completion to 98.4\% / 87.5\%. The filtered
radio reference reached 96.9\% / 93.8\%. A three step lookahead design initially
postponed movement at each replan, reaching 70.3\% / 53.1\%; applying progress
preference to the first predicted step yielded 93.8\% / 87.5\%, still with three
standard timeouts. These development failures are retained, not hidden by the
final controller selection.

In the second pilot design, the current station/group pair remained legal in
addition to explicit stay, eliminating an artificial alternating resource mask.
The third design introduced the residual braking prior and 200--1800 m training
range. No communication threshold was relaxed. The full final pilot already
showed fewer handovers and higher delay, motivating explicit component reporting
in the frozen comparison. These simultaneous development changes do not identify
individual causal contributions.

The following artifact families give a reproducible route from source data to
this paper. They are paths relative to the repository root; private binaries
are retained locally. Markdown notes 17--27 cover Study I and its audit, notes
28--32 cover Study II, and note 33 records this consolidation and publication.

{\footnotesize\setlength{\tabcolsep}{3pt}
\begin{longtable}{@{}p{0.16\textwidth}p{0.79\textwidth}@{}}
\caption{Artifacts and reproducibility boundaries.}\label{tab:artifacts}\\
\toprule Artifact & Location / purpose\\\midrule
\endfirsthead
\multicolumn{2}{l}{\textit{Continued from previous page}}\\\toprule
Artifact & Location / purpose\\\midrule
\endhead
\bottomrule\endfoot
Original map & dataset/Barcelona\_dataset\_January.h5; unchanged private HDF5, excluded from Git; SHA256 is stored in both freeze manifests.\\
Study I freeze & configs/frozen\_comparison\_v1.json; six exact source/protocol hashes.\\
Study II freeze & configs/frozen\_connectivity\_v2.json; seven exact source/protocol hashes.\\
Routes & configs/controlled\_scenarios.json and configs/connectivity\_scenarios\_v2.json; saved coordinate pairs, phases and seeds.\\
Outcomes & results/controlled\_experiment/confirmatory\_v1/ and results/connectivity\_experiment/confirmatory\_v2/; final evaluation records, training logs and source snapshots.\\
Statistics & Corresponding analysis\_v1/ and analysis\_v2/ directories; raw estimators, paired intervals, generated figures and table inputs. Study II report\_statistics.json documents the SINR reporting correction.\\
Replay & results/controlled\_experiment/connectivity\_audit\_v1/ and results/connectivity\_experiment/replay\_v2/; 4,400 audit and 5,600 final exact replays respectively.\\
Paper & paper/unified.tex and paper/unified\_appendices.tex; shared generated tables and figures; final consolidated IEEE format draft.\\
\end{longtable}}
